\newcommand{\protk}[1]{
    \prot^{(k)}(\rho(i_{#1 +1},\ldots,i_t = *),\rho(i_{#1} \to \bar{a_{#1}},i_{#1 +1},\ldots,i_t = *)) = \ts_#1
}

%%%%%%%%%%%%%%%%%%%%%
%E_k
\newcommand{\Et}{
    \substack{
        (i) \nvf(\phi_{a_1,\ldots,a_{t-1}}(\rho)) = i_t \hfill \\ 
%        (ii) \rho(i_j) = a_j \hfill\\
        (ii) F|_{\phi_{a_1,\ldots,a_{t}}(\rho)} \equiv 0, F|_{\phi_{a_1,\ldots,\bar{a_t}}(\rho)} \equiv 1 \hfill
        }
    }

\newcommand{\eve}[1]{
    \substack{
        \forall #1 \leq j \leq t \hfill\\
        (i) \nvf(\phi_{a_1,\ldots,a_{j-1}}(\rho)) = i_j \hfill \\ 
%        (ii) \rho(i_j) = a_j \hfill\\
        (ii) F|_{\phi_{a_1,\ldots,a_{t}}(\rho)} \equiv 0, F|_{\phi_{a_1,\ldots,\bar{a_t}}(\rho)} \equiv 1 \hfill
        }
    }
%%%%%%%%%%%%%%%%%%%%%%%%%

%%%%%%%%%%%%%%%%%%%%%
%Ft
\newcommand{\Ft}{
    \substack{
        (i)\nvf(\rho) = i_t \hfill \\ 
%        (ii) \rho(i_t) = * \hfill\\
        (ii) F|_{\rho(i_t \to a_t)} \equiv 0, F|_{\rho(i_t \to \bar{a_t})} \equiv 1 \hfill
        }
    }

% \newcommand{\Ft}{
%     \substack{
%         (i) \prot(\phi_{a_1,\ldots,a_{t-1}}(\rho)(i_t \to a_t),\phi_{a_1,\ldots,a_{t-1}}(\rho)(i_t \to \bar{a_t})) = \ts_t \hfill \\ 
% %        (ii) \rho(i_j) = a_j \hfill\\
%         (ii) F|_{\phi_{a_1,\ldots,a_{t-1}}(\rho)(i_t \to a_t)} \equiv 0, F|_{\phi_{a_1,\ldots,a_{t-1}}(\rho)(i_t \to \bar{a_t})} \equiv 1 \hfill
%         }
%     }
\newcommand{\Fkplusone}{
    \substack{
        \forall k+1 \le j \le t \hfill\\
        (i) \nvf(\phi_{a_{k,\ldots,a_{j-1}}}(\rho) ) = i_j \footnote{$\sigma_{k \to k-1}$ = ()}\hfill\\
        (ii) \rho(i_j) = * \hfill\\
        (iii) F|_{\phi_{a_k,\ldots,a_{t}}(\rho)} \equiv 0, F|_{\phi_{a_k,\ldots,\bar{a_t}}(\rho)} \equiv 1 \hfill
        }
    }
%%%%%%%%%%%%%%%%%%%%%%%%%    

%%%%%%%%%%%%%%%%%%%%%
%G_t
\newcommand{\Rlt}{
    \substack{
        (i) \prot(\rho(i_t \to a_t),\rho(i_t \to \bar{a_t})) = \ts_t \hfill \\ 
        (ii) \rho(i_t) = * \hfill\\
        (iii) F|_{\rho(i_t \to a_t)} \equiv 0, F|_{\rho(i_t \to \bar{a_t})} \equiv 1 \hfill
        }
    }
%%%%%%%%%%%%%%%%%%555

\newcommand{\Rlk}{
    \substack{
        (i) \prot^{(k)}(\rho(i_k \to a_k),\rho(i_k \to \bar{a_k})) = \ts_k \hfill \\ 
        (ii) \rho(i_k) = * \hfill\\
%        (iii) F|_{\rho(i_t \to a_t)} \equiv 0, F|_{\rho(i_t \to \bar{a_t})} \equiv 1 \hfill
        }
    }


%%%%%%%%%%%%%%%%%%%%%
\newcommand{\Gtprime}{
    \substack{
        (i) \prot(\rho(i_t\to a_t),\rho) = \ts_t \hfill \\ 
        (ii) \rho(i_t) = \bar{a_t} \hfill\\
        (iii) F|_{\rho(i_t \to a_t)} \equiv 0, F|_\rho \equiv 1 \hfill
        }
    }

\newcommand{\Gkprime}{
    \substack{
        \forall k \le j \le t \hfill\\
        (i) \prot(\rho(i_k\to a_k),\rho) = \ts_k \hfill \\ 
        (ii) \rho(i_k) = \bar{a_k} \hfill\\
        (iii) F|_{\rho(i_t \to a_t)} \equiv 0, F|_\rho \equiv 1 \hfill
        }
    }
%%%%%%%%%%%%%%%%%%%%%%%%%

%%%%%%%%%%%%%%%%%%%%%%%%%%%
% \newcommand{\evv}{
%     \substack{
%         \forall 1 \leq j \leq t \hfill\\
%         (i) \prot^{(j)}(\rho \circ \sigma_j,\rho\circ\sigma_j')= \ts_j \hfill\\ 
%         (ii) \rho(i_j) = * \hfill\\
%         (iii) F|_{\rho \circ\sigma_t} \equiv 0, F|_{\rho\circ \sigma_t'} \equiv 1 \hfill
%         }
%     }
%%%%%%%%%%%%%%%%%%%%%%%%%%%

%%%%%%%%%%%%%%%%%%%%%%%%%%
% \newcommand{\ev}[1]{
%     \substack{
%         \forall #1 \leq j \leq t \hfill\\
%         (i) \nvf(F,\phi_{a_{1},\ldots,a_{j-1}}(\rho)) = i_j \hfill \\ 
%         %(ii) \rho(i_j) = * \hfill\\
%         (ii) F|_{\phi_{a_1,\ldots,a_t}(\rho)} \equiv 0, F|_{\phi_{a_1,\ldots,\bar{a_t}}(\rho)} \equiv 1 \hfill
%         }
%     }
%%%%%%%%%%%%%%%%%%%%%%%%%


%%%%%%%%%%%%%%%%%%%
\newcommand{\even}[1]{
    \substack{
        \forall #1 \leq j \leq t \hfill\\
        (i) \nvf(F,\rho \circ \sigma_{k+1 \to j}) = i_j  \\ 
        (ii) \rho(i_j) = a_j \hfill\\
        (iii) F|_{\rho} \equiv 0, F|_{\rho(i_t \to a_t)} \equiv 1 \hfill
        }
    }  
%%%%%%%%%%%%%%%%%%%%%%%%



\begin{lemma} 
    \label{lem: main lemma}
    $$\prob{\rho}{\cdt(F,\rho) \mathrm{\ has\ depth\ } t} \leq \ans$$ 
\end{lemma}


% \begin{claim}
%     \label{label: NVF Game correspondence}
%     Let $\rho \in R_k$, 
%     $$\nvf(F,\rho) = x \implies \prot^{(k)}(\rho\circ (x \to a), \rho\circ(x \to \bar{a})) = x\ \mathrm{and}\ \rho(x) = *$$
% \end{claim}
    
\begin{proof} 
    If $\cdt(F,\rho)$ has depth $t$, then there exist $a \in \cube{t}$
    and variables $x_{i_1}, \ldots , x_{i_t}$ such that for all $k \in [t]$, $\nvf(F, \phi_{a_1,\ldots,a_{k-1}}(\rho) )= i_k$
    for all $k \in [t]$\footnote{assume $\phi_{a_0}(\rho) = \rho$}.

\begin{align*}
    \prob{\rho}{\cdt(F,\rho) \mathrm{\ has\ depth\ } t}   
    &\leq \sum_{x_{i_1}, \dots , x_{i_t}}  
            \prob{\rho}{\eve{1}} \\
\text{Let\ } E_k := \bra{\eve{k}}\\
    &=\sum_{x_{i_1}, \dots , x_{i_t}} \prob{\rho}{E_t}\prod_{k=t-1}^1 \prob{\rho}{E_k|E_{k+1}}\\
    &= \sum_{x_{i_t}} \pr{E_t} \sum_{x_{i_{t-1}}} \pr{E_{t-1}|E_t}\ldots  \sum_{x_{i_1}} \pr{E_1|E_2}
\end{align*}

The last step follows from the observation that for each $k \in [t]$, the event $E_k$ depends on the variables $x_{i_k}, \ldots, x_{i_t}$ and doesn't depend on the variables $x_{i_1}, \ldots , x_{i_{k-1}}$.

Let's begin with an upper bound the expression $\sum_{x_{i_t}} \pr{E_t}$. Then we move on to upper-bounds on the expressions $\sum_{x_{i_k}} \pr{E_k|E_{k+1}}$ which follow a similar line.

%%%%%%%%%%%%%%%%%%%%%%%%
%upperbound on E_t

Recall the definition of the event $E_t = \bra{\rho:\Et}$. 
\begin{align*}
    \sum_{x_{i_t}} \prob{\rho}{E_t}
        &= \sum_{x_{i_t}} \sum_{\rho \in E_t} \pr{\rho}\\
        &= \cbra{\frac{2p}{1-p}}^t \sum_{x_{i_t}} \sum_{\rho \in E_t} \pr{\phi_{a_1,\ldots, a_t}(\rho)}
\end{align*}        
        $\phi_{a_1,\ldots,a_t}(\rho)$ has $t$ more variables set 
        than $\rho$. Therefore $\pr{\rho}$ = 
        $\cbra{\frac{2p}{1-p}}^t$ 
        $\pr{\phi_{a_1,\ldots,a_t}(\rho)}$.\\
        Let $F_t:= \bra{\rho:\Ft}$.
\begin{align*}
    \sum_{\rho \in E_t} \pr{\phi_{a_1,\ldots, a_t}(\rho)}
    &\le  \sum_{\rho \in F_t} \pr{\rho(i_t \to a_t)} 
\end{align*}
 Let 
$R_{\ell_t} :=\bra{\rho:\Rlt}$. 
\begin{align*}
    \sum_{\rho \in F_t} \pr{\rho(i_t \to a_t)}
        &\leq \sum_{\ell_t} \sum_{\rho \in R_{\ell_t}} \pr{\rho(i_t \to a_t)}\\
        &= \sum_{\ell_t}\sum_{\rho \in R_{\ell_t}} \pr{\rho(i_t \to \bar{a_t})}\\
        &=\sum_{\ell_t} \sqrt{\sum_{\rho \in R_{\ell_t}} \pr{\rho(i_t \to a_t)}. \sum_{\rho \in R_{\ell_t}}\pr{\rho(i_t \to \bar{a_t})}}\\
        &\leq \sum_{\ell_i} \sqrt{ \pr{A_{\ts_t}} \pr{B_{\ts_t}} }
\end{align*}              


Let $U_t = \bra{\rho: F(\phi_{a_t}(\rho)) \equiv 0, F(\phi_{\bar{a_t}(\rho)} \equiv 1) }$.
Consider the communication game $\game_{F,U_t,a_t}$ as defined in sub-section \ref{sec: game}.
If $\rho \in U_t$ and $\nvf(\rho) = i_t$, the protocol $\prot^{(t)}$ recovers $i_t$ from $\phi_{a_t}(\rho)$ and $\phi_{\bar{a_t}}(\rho)$. 
Moreover, there exists a transcript $\ts_t$ of the protocol 
such that $\prot^{(t)}(\phi_{a_t}(\rho),\phi_{\bar{a_t}}(\rho)) = \ts_t $. Let $A_{\ts_t} \times B_{\ts_t}$
be the rectangle corresponding to the transcript $\ts_t$

%Let $\prot^{(t)}$ be the protocol for this game described in Lemma \ref{lem: protocol}.
Let $\prot^{(t)}$ be the Karchmer-Wigderson protocol for $F$. 
The protocol induces a partition of the communication matrix $M_{F,R_{t+1},a_t}$ into $L(F)$ many rectangles.
    Consider the following distributions $p_t$ and $q_t$ supported on $X_t = \{\phi_{a_t}(\rho) : \rho \in U_t\}$ and $Y_t = \{\phi_{\bar{a_t}}(\rho): \rho \in U_t\}$ respectively.
    For all $\rho \in X_t$, 
            $$p_t[\rho] = \prob{\rho \in R_p}{\rho},$$
        and for all $\rho \in Y_t$, 
            $$q_t[\rho] = \prob{\rho \in R_p}{\rho}.$$ 
        Note that $\sum_{\rho \in X_t} p_t[\rho] \leq 1$ and similarly 
    $\sum_{\rho \in Y_t} q_t[\rho] \leq 1$.
    Let 
    $p_t[A_{\ell_t}] := \sum_{\rho \in R_{\ell_t}} \pr{\phi_{a_t}(\rho)}$ and $p_t[B_{\ell_t}] := \sum_{\rho \in R_{\ell_t}} \pr{\phi_{\bar{a_t}}(\rho)}$.
    By Lemma \ref{lem: sum of rectangles},
    $\sum_{x_{i_t}}\sum_{\ts_t \in \prot^{(t)}} p_t[A_{\ts_t}]q_t[B_{\ts_t}] \leq 1$.

\begin{align*}
    \sum_{x_{i_t}}\sum_{\ts_t} \sqrt{ p_{t}[A_{\ts_t}] q_t
    [B_{\ts_t}] }
        &\leq \sqrt{\sum_{x_{i_t},\ts_t} 1} \sqrt{\sum_{x_{i_t},\ts_t}p_t[A_{\ts_t}] q_t[B_{\ts_t}]}\\
        &\leq \sqrt{|\prot^{(t)}|} \leq \sqrt{L(F)}.
\end{align*}

Since $\prot^{(t)}$ is the KW protocol which has $L(F)$ many transcripts. 
%%%%%%%%%%%%%%

%%%%%%%%%%%
%upperbound on E_k|E_k+1
We now upper bound $\sum_{x_{i_k}} \prob{\rho}{E_k|E_{k+1}}$.
Recall the definition of the event $E_k = \bra{\eve{k}}$.
\begin{align*}
    \sum_{x_{i_k}} \prob{\rho}{E_k|E_{k+1}}
    &= \sum_{x_{i_k}} \frac{\sum_{\rho \in E_k} \pr{\rho}}{\sum_{\rho \in E_{k+1}} \pr{\rho}}\\
            &= \sum_{x_{i_k}} \frac{\sum_{\substack{\rho \in E_{k+1}\\ \nvf(\phi_{a_1,\ldots, a_{k-1}}(\rho)) = i_k}} \pr{\phi_{a_1,\ldots, a_k}(\rho)}}{\sum_{\rho \in E_{k+1}} \pr{\phi_{a_1,\ldots, a_k}(\rho)}}\\
\text{Let}\ F_{k+1}:= \bra{\rho:\Fkplusone}\\
    &\leq \sum_{x_{i_k}} \frac{\sum_{\substack{\rho \in F_{k+1}\\ \nvf(\rho) = i_k}} \pr{\phi_{a_k}(\rho)}}{\sum_{\rho \in F_{k+1}} \pr{\phi_{a_k}(\rho)}}\\
\end{align*}
 Let 
$R_{\ell_k} :=\bra{\rho:\Rlk}$. 
\begin{align*}
    \frac{\sum_{\substack{\rho \in F_{k+1}\\ \nvf(\rho) = i_k}} \pr{\phi_{
    a_k}(\rho)}}{\sum_{\rho \in F_{k+1}} \pr{\phi_{a_k}(\rho)}}
%            &\leq \sum_{\ell_k} \sum_{\rho \in R_{\ell_k} \cap F_{k+1}} \pr{\rho(i_t \to a_t)}\\
%        &= \sum_{\ell_t}\sum_{\rho \in R_{\ell_t}} \pr{\rho(i_t \to \bar{a_t})}\\
        &=\sum_{\ell_t} 
            \sqrt{
                \frac{
                    \sum_{\rho \in R_{\ell_k} \cap F_{k+1}} 
                        \pr{\phi_{a_k}(\rho)}
                }
                {
                \sum_{\rho \in F_{k+1}} 
                    \pr{\phi_{a_k}(\rho)}
                }
            }
            \sqrt{
                \frac{
                    \sum_{\rho \in R_{\ell_k} \cap F_{k+1}} 
                        \pr{\phi_{\bar{a_k}}(\rho)}
                }
                {
                \sum_{\rho \in F_{k+1}} 
                    \pr{\phi_{\bar{a_k}}(\rho)}
                }
            }\\
        &\leq \sum_{\ell_i} \sqrt{ \pr{A_{\ts_k}} \pr{B_{\ts_k}} }
\end{align*}            

Consider the communication game $\game_{F,F_{t+1},a_k}$ as defined in sub-section \ref{sec: game}.
Let $\prot^{(k)}$ be the protocol for this game described in Lemma \ref{lem: protocol}.
If $\rho \in F_{k+1}$ and $\nvf(\rho) = i_k$, the protocol $\prot^{(k)}$ recovers $i_k$ from $\phi_{a_k}(\rho)$ and $\phi_{\bar{a_k}}(\rho)$. 
Moreover, there exists a transcript $\ts_k$ of the protocol 
such that $\prot^{(k)}(\phi_{a_k}(\rho),\phi_{\bar{a_k}}(\rho)) = \ts_k $. Let $A_{\ts_k} \times B_{\ts_k}$
be the rectangle corresponding to the transcript $\ts_k$. 


The protocol induces a partition of the communication matrix $M_{F,F_{k+1},a_k}$ into $L(F)^2$ many rectangles.
    Consider the following distributions $p_k$ and $q_k$ supported on $X_k = \bra{\phi_{a_k}(\rho):\rho \in F_{k+1}}$ and $Y_k= \bra{\phi_{\bar{a_k}}(\rho): \rho \in F_{k+1}}$ respectively.
    For all $\rho \in X_t$, 
            $$p_k[\phi_{a_k}(\rho)] = \frac{
                                \pr{\phi_{a_k}(\rho)}
                            }
                            {
                            \sum_{\rho \in F_{k+1}}
                                \pr{\phi_{a_k}(\rho)}       
                            },$$
        and for all $\rho \in Y_t$, 
            $$q_k[\phi_{\bar{a_k}}(\rho)] = \frac{
                                \pr{\phi_{\bar{a_k}}(\rho)}
                            }
                            {
                            \sum_{\rho \in F_{k+1}}
                                \pr{\phi_{\bar{a_k}}(\rho)}               
                            },$$

        Note that $p_k$ and $q_k$ are distributions on $X_k$ and $Y_k$ respectively i.e. $\sum_{\rho \in X_k} p_k[\rho] = 1$ and $\sum_{\rho \in Y_k} q_k[\rho] = 1$.
    Let $R_{\ell_k} = A_{\ell_k} \times B_{\ell_k}$ where $A_{\ell_k} \subseteq X_k$ and $B_{\ell_k} \subseteq Y_k$. $p_k[A_{\ell_k}] = \sum_{\rho \in R_{\ell_k} \cap F_{k+1}} p_k[\phi_{a_k}(\rho)] $ and $q_k[B_{\ell_k}] = \sum_{\rho \in R_{\ell_k} \cap F_{k+1}} q_k[\phi_{\bar{a_k}}(\rho)]$. 
    By Lemma \ref{lem: sum of rectangles},
    $\sum_{x_{i_k}}\sum_{\ts_k \in \prot^{(k)}} p_k[A_{\ts_k}]q_k[B_{\ts_k}] = 1$. Therefore

\begin{align*}
    \sum_{x_{i_k}}\sum_{\ts_k} \sqrt{ p_{k}[A_{\ts_k}] q_k
    [B_{\ts_k}] }
        &\leq \sqrt{\sum_{\ts_k} 1} \tag*{by Cauchy-Schwarz} \sqrt{\sum_{\ts_k}p_k[A_{\ts_k}] q_k[B_{\ts_k}]}\\
        &\leq \sqrt{|\prot^{(k)}|} \leq L(F).
\end{align*}
    \begin{align*}
        \prob{\rho}{\cdt(F,\rho) \mathrm{\ has\ depth\ } t} 
            &= \sum_{x_{i_t}} \pr{E_t} \sum_{x_{i_{t-1}}} \pr{E_{t-1}|E_t}\ldots  \sum_{x_{i_1}} \pr{E_1|E_2}\\
            &\leq p^t \sqrt{L(F)} L(F)^{t-1}            
    \end{align*}
\end{proof}



%%%%%%%%%%%%%%%%%%%%%%%%%%%%%%%%%%%%%%%%%%%%%%%%%%%%%
        % The above equality follows from the following observations 
        % \begin{itemize}
        %     \item $\rho$ can be recovered from $\rho \circ (i_1 \to a_1, \ldots , i_t \to a_t)$ given $(i_1,\ldots ,i_t)$ and $a \in \cube{t}$.
        %     \item $\frac{\prob{}{\rho}}{p^t}  = \frac{\prob{}{\rho\circ \sigma_t}}{((1-p)/2)^t} $ when $\rho(i_1) = \ldots = \rho(i_t)= *$.
        % \end{itemize} 
          
        % %%%%%%%%%%%%  
        % \newcommand{\prottk}[1]{\prot^{(k)}(\rho(i_{#1} \to a_{#1}),\rho(i_{#1} \to \bar{a_{#1}})) = \ell_{#1}}
        % %%%%%%%%%%%%%%%%
        
        % \begin{align*}
        %     % \prob{\rho}{(i),(ii) \ \text{hold for}\ 1\leq k \leq t,(iii)}
        %     % &= \prod_{k=t}^1 \prob{\rho}{(i),(ii)\ \text{hold for}\ k \leq j \leq t, (iii)\bigg| (i),(ii)\ \text{hold for}\ k+1 \leq j \leq t,(iii)}\\
        %     &\prob{\rho}{\eve{1}}\\
        %     &= \prod_{k=t}^1 \prob{\rho}{\eve{k} \middle| \eve{k+1}}
        %     \footnote{for k= t, the even we are conditioning is vaccuously hold for all $\rho$. We do that only for notational convenience.}\\
        %     \tag{factorizing the event into conditional events}\\
        % %
        %     &= \prod_{k=t}^1 
        %     \prob{\rho}{
        %         \substack{
        %             \rho \in E_{k+1}\\
        %             \rho(i_k) = a_k \\ 
        %             \protk{k}
        %             %\nvf(F,\rho(i_k,\ldots,i_t \to *)) = i_k
        %         }
        %     }\bigg/\prob{\rho}{\rho \in E_{k+1}}\\
        %     \tag{$E_k$ is shorthand for the event}\\
        % %
        %     &= \prod_{k=t}^1 
        %     \frac{
        %         \sum_{
        %             \substack{
        %                 \rho \in F_{k+1}, \rho(i_k) = a_k\\ 
        %                 \prot^{(k)}(\rho,\rho(i_k \to \bar{a_k})) = \ts_k
        %                 %\nvf(F,\rho)=i_k 
        %             }
        %         } \pr{\rho}
        %     }
        %     {\sum_{
        %         \rho \in F_{k+1}
        %         } \pr{\rho}
        %     }\\
        %     \tag{substituting $E_k$ $\rho(i_{k+1},\ldots,i_t = *)$ with $\rho$ $F_k$}\\
        % %   
        %     &\leq \prod_{k=t}^1 
        %         \sqrt{
        %             \frac{
        %                 \sum_{
        %                     \substack{
        %                         \rho \in F_{k+1}, \rho(i_k)= a_k\\
        %                         \prot^{(k)}(\rho,\rho(i_k \to \bar{a_k})) = \ts_k
        %                         %\prot^{(k)}(\rho,\rho(i_k \to \bar{a_k})) = \ts_k 
        %                         }
        %                     } \pr{\rho}
        %                 }
        %                 {\sum_{
        %                     \rho \in F_{k+1}
        %                     } \pr{\rho}
        %                 }
        %             }
        %         \sqrt{
        %             \frac{
        %                 \sum_{
        %                     \substack{
        %                         \rho \in F_{k+1}, \rho(i_k= *)\\
        %                         \prot^{(k)}(\rho(i_k \to a_k),\rho) = \ts_k
        %                         %\prot^{(k)}(\rho(i_k \to a_k),\rho) = \ts_k
        %                         }
        %                     } \pr{\rho}
        %                 }
        %                 {\sum_{
        %                     \rho \in F_{k+1}
        %                     } \pr{\rho}
        %                 }
        %             }\\
        %     &\tulasi{this is a crutial step. Check.}\\
        % %    
        %     &=  \sqrt{
        %             \frac{
        %                 \prob{}{A_k(\ts_k)}
        %             }
        %             {
        %                 \pr{X_k}
        %             }
        %         }
        %         \sqrt{
        %             \frac{
        %                 \pr{B_k(\ts_k)}
        %             }
        %             {
        %             \prob{}{Y_k}
        %             }
        %         } 
        %     % &= \prod_{k=t}^1 \prob{\rho}{ \substack{\rho \in F_{k+1}\\
        %     % \rho(i_k) = a_k \\ \nvf(F,\rho(i_k \to *)) = i_k
        %     % }\middle| \rho \in F_{k+1}}\\
        %     % &= \prod_{k=t}^1 \frac{\sum_{\rho \in \calA_k} \prob{}{\rho}}{\sum_{\rho \in \calA_{k+1}} \prob{}{\rho}} 
        % \end{align*}
        
        % Where, $E_k= \{
        %         \rho :
        %         \protk{j}
        %         %\nvf(F,\rho(i_j,\ldots,i_t \to *)) = i_j
        %         , \rho(i_j) = a_j \forall k \leq j \leq t 
        %     \}$
            
        %     % $E_{k+1}= \{
        %     %     \rho: 
        %     %     \protk{j}
        %     %     %\nvf(F,\rho(i_j, \ldots , i_t \to *)) = i_j 
        %     %     , \rho(i_j) = a_j  \forall k+1 \leq j \leq t   
        %     % \}$
        
        %     % $E_k = \{ 
        %     %     \rho \in E_{k+1}:\ \rho(i_k) = a_k,\ 
        %     %     \protk{k}
        %     %     %\nvf(F,\rho(i_k, \ldots , i_t \to *)) = i_k 
        %     % \}$ 
            
        %     $F_{k+1} =\{ 
        %         \rho:
        %         \prot^{(j)}(\rho\circ \sigma_{k+1 \to j}, \rho \circ \sigma_{k+1 \to j}')
        %         %\nvf(F, \rho(\sigma_{k \to j-1}))= i_j
        %         , \rho(i_j) = *\ \forall k+1 \leq j \leq t      
        %     \}$
        
        % % where $\calA_k : = \{
        % %     \rho: (i),(ii)\ \text{hold for}\ k \leq j \leq t, (iii) 
        % % \}$
        
        % %$X_{t+1} = \text{All Restrictions}$.
        
        % $X^{(t)} = \{
        %     \phi_{a_t}(\rho) :  F|_{\phi_{a_t}(\rho)} \equiv 0, F|_{\phi_{\bar{a_t}}}(\rho) \equiv 1  
        % \}$
                
        % $Y^{(t)} = \{
        %     \phi_{\bar{a_t}}(\rho) : \rho \in X_{t+1}, F|_{\phi_{a_t}}(\rho) \equiv 0, F|_{\phi_{\bar{a_t}}(\rho)}\equiv 1  
        % \}$        
            
        % % $X_{t-1} = \{
        % %         \phi_{a_{t-1}}(\rho): \rho \in X_t, \rho(i_t) = *, \prot^{(t)}(\rho(i_t \to a_t),\rho(i_t \to \bar{a_t})) = \ts_t     
        % %         \}$        
        
        % % $Y_{t-1}= \{
        % %         \phi_{\bar{a_{t-1}}}(\rho): \rho \in X_t, \rho(i_t) = *,
        % %         \prot^{(t)}(\rho(i_t \to a_t),\rho(i_t \to \bar{a_t})) = \ts_t
        % %         \}$
        
        % % $ X_k = \{
        % %         \phi_{a_k}(\rho): \rho \in X_{k+1}, \rho(i_{k+1}) = *,
        % %         \prot^{(k+1)}(\rho(i_{k+1} \to a_{k+1}),\rho(i_{k+1} \to \bar{a_{k+1}})) = \ts_{k+1}
        % %         \}$ keept his 
        
        % $X^{(k)} = \{
        %         \phi_{a_k}(\rho): \rho \in \calA_{k+1} 
        % \}$
        
        % $Y^{(k)} = \{
        %         \phi_{\bar{a_k}}(\rho): \rho \in \calA_{k+1} 
        % \}$
        
                
        % % $ Y_k = \{
        % %         \phi_{\bar{a_k}}(\rho): \rho \in X_{k+1}, \rho(i_{k+1}) = *,
        % %         \prot^{(k+1)}(\rho(i_{k+1} \to a_{k+1}),\rho(i_{k+1} \to \bar{a_{k+1}})) = \ts_{k+1}
        % %         \}$ keep this         
                
        % $A^{(k)}(\ell_k) = \{
        %         \rho: \rho \in X^{(k)}, \prot^{(k)}(\rho, \rho(i_k \to \bar{a_k})) = \ts_k
        %         \}$
        
        % $B^{(k)}(\ell_k) = \{
        %         \rho: \rho \in Y^{(k)}, \prot^{(k)}(\rho(i_k \to a_k), \rho) = \ts_k
        % \}$
                
        % %  \sum_{x_{i_1}, \ldots , x_{i_t}}
        % %      \sum_{\rho \in \calA} Pr[\rho]
        
        % % Fix a choice of $x_{i_1},\ldots , x_{i_t}$ and $a \in \cube{t}$. For each $j \in [t]$, let 
        % % $\sigma_j := \sigmaj{j}$ 
        % % and $\sigma_j' := \sigmaprimej{j}$. Let $\sigma_0 = \sigma_0' = ()$.
        
        
        % \begin{align*}
        %     \sum_{x_{i_k}} 
        %         \sum_{\ts_k}    
        %             \sqrt{
        %                 \frac{
        %                     \pr{A_k(\ts_k)}
        %                 }
        %                 {\pr{X_k}}
        %             }
        %             \sqrt{
        %                 \frac{
        %                     \pr{B_k(\ts_k)}
        %                 }
        %                 {\pr{Y_k}}
        %             } 
        %     &\leq \sqrt{|\prot^{(k)}|}        
        %     \sqrt{
        %         \sum_{x_k,\pi_k}
        %             \frac{
        %                 \pr{A_k(\ts_k)}
        %             }
        %             {
        %             \pr{X_k}
        %             }
        %             \frac{
        %                 \pr{B_k(\ts_k)}
        %             }
        %             {\pr{Y_k}}
        %     }\\
        %     &= \sqrt{|\prot^{(k)}|} 
        % \end{align*}
                
        %     Therefore, $\prob{\rho}{\cdt(F,\rho) \mathrm{\ has\ depth\ } t} \leq p^t \prod_{k=1}^t \sqrt{|\prot^{(k)}|}$
        
        
        % %\newcommand{\A}{\left\{
    \rho: \prot(\rho\circ \sigma_j, \rho \circ \sigma_j')= \ts_j,\ 
    \rho(i_j) = *\ \forall 1 \leq j \leq t;\
    F|_{\rho\circ \sigma_t} \equiv 0,\
    F|_{\rho \circ \sigma_t'} \equiv 1  
    \right\}}

\newcommand{\Bk}[1]{\{
    \rho: \prot(\rho(i_{j+1} = \ldots = i_t = *) ), \rho(i_j \to \bar{a_j}, i_{j+1} = \ldots = i_t = *))= \ts_j,\ 
    \rho(i_j) = a_j\ \forall #1 \leq j \leq t;\
    F|_{\rho} \equiv 0,\
    F|_{\rho(i_t \to \bar{a_t})} \equiv 1  
    \}}

\newcommand{\Ck}[1]{\{
    \rho: \prot(\rho, \rho(i_k \to \bar{a_k})= \ts_k;\prot(\rho\circ \sigma_{k+1 \to j}, \rho(i_k \to \bar{a_k})\circ \sigma_{k+1 \to j}')= \ts_j, \rho(i_j) = * \forall #1+1 \leq j \leq t ;
    \rho(i_k) = a_k ,
    F|_{\rho\circ \sigma_{k+1 \to t}} \equiv 0,
    F|_{\rho\circ \sigma_{k+1 \to t}'} \equiv 1  
    \}}
\newcommand{\Ckprime}[1]{\{
    \rho: \prot(\rho(i_k \to a_k), \rho)= \ts_k;\prot(\rho(i_k \to a_k)\circ \sigma_{k+1 \to j}, \rho(i_k \to a_k)\circ \sigma_{k+1 \to j}')= \ts_j, \forall #1+1 \leq j \leq t ;
    \rho(i_k) = \bar{a_k} ,
    F|_{\rho\circ(i_k \to a_k) \sigma_{k+1 \to t}} \equiv 0,
    F|_{\rho\circ(i_k \to a_k) \sigma_{k+1 \to t}'} \equiv 1 
    \}}
    
\newcommand{\calD}{\mathcal{D}}

\newcommand{\Dk}[1]{\{
    \rho: \prot(\rho(i_j \to a_j)\circ \sigma_{k+1 \to j}, \rho\circ \sigma_{k+1 \to j}')= \ts_j, \forall #1 \leq j \leq t ;
    \rho(i_k) = a_k ,
    F|_{\rho\circ \sigma_{k+1 \to t}} \equiv 0,
    F|_{\rho(i_j \to \bar{a_j})\circ \sigma_{k+1 \to t}'} \equiv 1  
    \}}
\newcommand{\Dkprime}[1]{\{
    \rho: \prot(\rho\circ \sigma_{k+1 \to j}, \rho(i_j \to \bar{a_j})\circ \sigma_{k+1 \to j}')= \ts_j, \forall #1+1 \leq j \leq t ;
    \rho(i_k) = a_k ,
    F|_{\rho\circ \sigma_{k+1 \to t}} \equiv 0,
    F|_{\rho(i_j \to \bar{a_j})\circ \sigma_{k+1 \to t}'} \equiv 1  
    \}}


where $\calA := \A$.
%wite \calA interms of X_k , Y_k A_k B_k


\begin{align*}
        \sum_{\rho \in \calA} \prob{}{\rho} 
        &= \cbra{\frac{2p}{1-p}}^t \sum_{\rho \in \calA} \prob{}{\rho \circ \sigma_t}\\
        &= \cbra{\frac{2p}{1-p}}^t  \sum_{\rho \in \calB} \prob{}{\rho} 
        \\
\end{align*}        

Where $\calB:= 
    \left\{ \rho:
        F|_\rho \equiv 0, F|_{\rho(i_t \to a_t)} \equiv 1;  
        \rho(i_k) = a_k,\
        \nvf(F,\rho(i_k, \ldots ,i_t \to *)) = i_k\ \forall 1 \leq k \leq t;        
    \right\}
$

Consider the communication games $G_1, \dots , G_t$ 

Where $G_k = \game_{F,\calR_k,a_k}$. 

For example 
$\calR_t$




where $\calB(i_1,\ldots ,i_k) = \Bk{1}$.

\begin{align*}
        \sum_{x_{i_1}, \ldots , x_{i_t}}   \sum_{\substack{\ts_1 \in \mathcal{L}_1\\ \vdots \\ \ts_t \in \mathcal{L}_t}}\sum_{\rho \in \calB} \prob{}{\rho}
        &= \sum_{x_{i_1}, \ldots , x_{i_t}}\sum_{\substack{\ts_1 \in \mathcal{L}_1\\ \vdots \\ \ts_t \in \mathcal{L}_t}}\sum_{\rho \in \calB_t} \prob{}{\rho} \cdot \prod_{k =1}^{t-1} \frac{\sum_{\rho \in \calB_k} \prob{}{\rho}}{\sum_{\rho \in \calB_{k+1}} \prob{}{\rho}}\\
        &= \sum_{\substack{x_{i_t}\\ \ts_t \in \mathcal{L}_t}} S_t \sum_{\substack{x_{i_{t-1}}\\ \ts_{t-1} \in \mathcal{L}_{t-1}}} S_{t-1} \ldots \sum_{\substack{x_{i_1}\\ \ts_1 \in \mathcal{L}_1}} S_1
\end{align*}
\todo{it is not $B_{k+1}$}
where $\calB_k(i_k, \ldots , i_t) = \Bk{k}$, 
$S_k(i_k, \ldots , i_t) = \frac{\sum_{\rho \in \calB_k(i_k, \ldots , i_t)} \prob{}{\rho}}{\sum_{\rho \in \calB_{k+1}(i_{k+1}, \ldots ,i_t)}\prob{}{\rho}} \forall 1 \leq k \leq t-1$ and $S_t(i_t) = \sum_{\rho \in \calB_t(i_t)} \prob{}{\rho}$.\\
Note that $S_k$ depends on the variables $x_{i_t}, \ldots, x_{i_k}$. 

Consider the expression, 
    \begin{align*}
        \sum_{\substack{x_{i_k}\\ \ts_k \in \mathcal{L}_k }}S_k(i_t ,\ldots, i_k) 
        &= \sum_{\substack{x_{i_k}\\ \ts_k \in \mathcal{L}_k }} \frac{\sum_{\rho \in \calB_k(i_k ,\ldots, i_t) \prob{}{\rho}}}{\sum_{\rho \in \calB_{k+1}(i_{k+1} ,\ldots, i_t)}\prob{}{\rho}}\\
        &= \sum_{\substack{x_{i_k}\\ \ts_k \in \mathcal{L}_k }} \frac{\sum_{\rho \in \calC_k(i_k ,\ldots, i_t)} \prob{}{\rho}}{\sum_{\rho \in \calD_{k+1}(i_{k+1} ,\ldots, i_t)}\prob{}{\rho}}\\
        &= \sum_{\substack{x_{i_k}\\ \ts_k \in \mathcal{L}_k }} \frac{\sum_{\rho \in \calC_k(i_k ,\ldots, i_t)} \sqrt{\prob{}{\rho}}\sqrt{\prob{}{\rho \circ (x_{i_k \to \bar{a_k})}}}}{\sum_{\rho \in \calD_{k+1}(i_{k+1} ,\ldots, i_t)}\prob{}{\rho}}\\
        &\leq \sum_{\substack{x_{i_k}\\ \ts_k \in \mathcal{L}_k }} \frac{ \sqrt{\sum_{\rho \in \calC_k(i_k ,\ldots, i_t)}\prob{}{\rho}}\sqrt{\sum_{\rho \in \calC_k(i_k ,\ldots, i_t)}\prob{}{\rho \circ (x_{i_k \to \bar{a_k})}}}}{\sum_{\rho \in \calD_{k+1}(i_{k+1} ,\ldots, i_t)}\prob{}{\rho}}\\
        &= \sum_{\substack{x_{i_k}\\ \ts_k \in \mathcal{L}_k }} 
        \sqrt{\frac{\sum_{\rho \in \calC_k}\prob{}{\rho}}{{\sum_{\rho \in \calD_{k+1}}\prob{}{\rho}}}} \sqrt{\frac{\sum_{\rho \in \calC_k'}\prob{}{\rho}}{\sum_{\rho \in \calD_{k+1}}\prob{}{\rho}}}  
    \end{align*}
 
where $\calC_k = \Ck{k}$,\\ 
$\calC_k' = \Ckprime{k}$

where $\calD_{k+1} = \Dk{k+1}$

% \begin{align*}
% \frac{\sum_{\lambda \in \calB_k} \prob{}{\lambda}}{\sum_{\lambda \in \calB_{k+1}} \prob{}{\lambda}}
%         = \frac{\sum_{\eta \in \calC_k} \prob{}{\eta}}{\sum_{\eta \in \calC_{k+1}} \prob{}{\eta}} \tag{$\eta = \lambda(i_{k+1}= \ldots = i_{t} = *)$}  
%     \end{align*}

% $\calB_k' = \Bk{k}$


% %$\calB_k' = \Bkprime{k}$\\

% $\calC_k' = \Ckprime{k}$\\

% $\calD_k'= \Dkprime{k}$\\

% $\calD_k = \Dk{k}$\\

    


% $\sum_{\eta \in \calC_t} \prob{}{\eta} = \sqrt{\sum_{\eta \in \calC_t} \prob{}{\eta}}\sqrt{\sum_{\eta \in \calD_t} \prob{}{\eta}}$\\

% $\frac{\sum_{\eta \in \calC_k} \prob{}{\eta}}{\sum_{\eta \in \calC_{k+1}} \prob{}{\eta}} = \sqrt{\frac{\sum_{\eta \in \calD_k} \prob{}{\eta}}{\sum_{\eta \in \calD_k'} \prob{}{\eta}}} \cdot \sqrt{\frac{\sum_{\eta \in \calC_k} \prob{}{\eta}}{\sum_{\eta \in \calC_k'} \prob{}{\eta}}} $



%     \begin{align*}
%         \sum_{\rho \in \calA_1} Pr[\rho]
%         &= \sum_{\rho \in \calA_t} \prob{}{\rho} \cdot \prod_{k=1}^{t-1} \frac{\sum_{\rho \in \calA_k}\prob{}{\rho}}{\sum_{\rho \in \calA_{k+1}}\prob{}{\rho}}\\
%         &= \cdot \prod_{k=1}^{t-1} \frac{\sum_{\rho \in \calA_k}\prob{}{\rho}}{\sum_{\rho \in \calA_{k+1}}\prob{}{\rho}}\\
%     \end{align*}
   
%   \begin{align*}
%         \sum_{\rho \in \calA_t} \prob{}{\rho} 
%         &= \sum_{\rho \in \calA_t} \sqrt{\prob{}{\rho \circ \sigma_t}}\sqrt{\prob{}{\rho \circ \sigma_t'}}  \\
%         &\leq \sqrt{\sum_{\rho \in \calA_t}\prob{}{\rho \circ \sigma_t}}\sqrt{\sum_{\rho \in \calA_t}\prob{}{\rho \circ \sigma_t'}}
%   \end{align*}

%     \begin{align*}
%         \frac{\sum_{\rho \in \calA_k}\prob{}{\rho}}{\sum_{\rho \in \calA_{k+1}}\prob{}{\rho}} 
%         &= \sum_{\rho \in \calA_k} \sqrt{\frac{\prob{}{\rho\circ \sigma_k}}{\sum_{\rho \in \calA_{k+1}}\prob{}{\rho\circ \sigma_k}}} \sqrt{\frac{\prob{}{\rho\circ \sigma_k'}}{\sum_{\rho \in \calA_{k+1}}\prob{}{\rho\circ \sigma_k'}}} \\
%         &\leq \sqrt{\frac{\sum_{\rho \in \calA_k}\prob{}{\rho\circ \sigma_k}}{\sum_{\rho \in \calA_{k+1}}\prob{}{\rho\circ \sigma_k}}} \sqrt{\frac{\sum_{\rho \in \calA_k}\prob{}{\rho\circ \sigma_k'}}{\sum_{\rho \in \calA_{k+1}}\prob{}{\rho\circ \sigma_k'}}}
%     \end{align*}

    
%     $\calB_{k}:= \{
%     \rho: \prot(\rho, \rho(i_k \to \bar{a_k}))= \ts_k;
%     \prot(\rho \circ \sigma_{k+1 \to j}, \rho \circ \sigma_{k+1 \to j}')= \ts_j, \rho(i_j) = * \forall k+1 \leq j \leq t  ; 
%     \rho(i_j) = a_j \forall 1 \leq j \leq k;
%     F|_{\rho\circ \sigma_t} \equiv 0,
%     F|_{\rho \circ \sigma_t'} \equiv 1  
%      \}$
    
%     $\calB_{k}':= \left\{
%     \rho: \prot(\rho \circ \sigma_{k+1 \to j}, \rho \circ \sigma_{k+1 \to j}')= \ts_j, \rho(i_j) = * \forall k+1 \leq j \leq t  ; 
%     \rho(i_j) = a_j \forall 1 \leq j \leq k;
%     F|_{\rho\circ \sigma_t} \equiv 0,
%     F|_{\rho \circ \sigma_t'} \equiv 1  
%      \right\}$
    
%     $\calC_{k}:= \left\{
%     \rho: \prot(\rho(i_k \to a_k),\rho)= \ts_k;\prot(\rho(i_k \to a_k) \circ \sigma_{k+1 \to j}, \rho(i_k \to a_k) \circ \sigma_{k+1 \to j}')= \ts_j, \rho(i_j) = * \forall k+1 \leq j \leq t  ;
%     \rho(i_j) = a_j \forall 1 \leq j \leq k-1; \rho(i_k) = \bar{a_k};
%     F|_{\rho\circ \sigma_t} \equiv 0,
%     F|_{\rho \circ \sigma_t'} \equiv 1  
%      \right\}$
    
%     $\calC_{k}':= \left\{
%     \rho: \prot(\rho(i_k \to a_k) \circ \sigma_{k+1 \to j}, \rho(i_k \to a_k) \circ \sigma_{k+1 \to j}')= \ts_j, \rho(i_j) = * \forall k+1 \leq j \leq t  ;
%     \rho(i_j) = a_j \forall 1 \leq j \leq k-1;\\
%     \rho(i_k) = \bar{a_k};
%     F|_{\rho\circ \sigma_t} \equiv 0,
%     F|_{\rho \circ \sigma_t'} \equiv 1  
%      \right\}$
     
%     \begin{align*}
%         \frac{\sum_{\rho \in \calA_k}\prob{}{\rho\circ \sigma_k}}{\sum_{\rho \in \calA_{k+1}}\prob{}{\rho\circ \sigma_k}}
%         &= \frac{\sum_{\rho \in \calB_k} \prob{}{\rho}}{\sum_{\rho \in \calB_k'} \prob{}{\rho}}\\
%         &= \frac{\sum_{\substack{\rho:\rho \in \calB_k'\\ \prot(\rho, \rho(i_k \to \bar{a_k}))= \ts_k }} \prob{}{\rho}}{\sum_{\rho \in \calB_k'} \prob{}{\rho}}
%     \end{align*}
%     Similarly,
%     \begin{align*}
%         \frac{\sum_{\rho \in \calA_k}\prob{}{\rho\circ \sigma_k'}}{\sum_{\rho \in \calA_{k+1}}\prob{}{\rho\circ \sigma_k'}}
%         &= \frac{\sum_{\rho \in \calC_k} \prob{}{\rho}}{\sum_{\rho \in \calC_k'} \prob{}{\rho}}\\
%         &= \frac{\sum_{\substack{\rho:\rho \in \calC_k'\\ \prot(\rho(i_k \to a_k), \rho)= \ts_k }} \prob{}{\rho}}{\sum_{\rho \in \calC_k'} \prob{}{\rho}}
%     \end{align*}    
        
        
        % %\newcommand{\A}{\left\{
    \rho: \nvf(\rho \circ \sigma_{k-1} = x_{i_k}) \forall1 \leq  k \leq t, \rho(i_j) = * \forall  1\leq j \leq t, F|_{\rho \circ \sigma_t \equiv 0} , F|_{\rho \circ \sigma_t'} \equiv 1
    \right\}}

\newcommand{\Bk}[1]{\{
    \rho: \ts(\rho(i_{j+1} = \ldots = i_t = *) ), \rho(i_j \to \bar{a_j}, i_{j+1} = \ldots = i_t = *))= \ts_j,\ 
    \rho(i_j) = a_j\ \forall #1 \leq j \leq t;\
    F|_{\rho} \equiv 0,\
    F|_{\rho(i_t \to \bar{a_t})} \equiv 1  
    \}}

\newcommand{\Ck}[1]{\{
    \rho: \ts(\rho, \rho(i_k \to \bar{a_k})= \ts_k;\ts(\rho\circ \sigma_{k+1 \to j}, \rho(i_k \to \bar{a_k})\circ \sigma_{k+1 \to j}')= \ts_j, \rho(i_j) = * \forall #1+1 \leq j \leq t ;
    \rho(i_k) = a_k ,
    F|_{\rho\circ \sigma_{k+1 \to t}} \equiv 0,
    F|_{\rho\circ \sigma_{k+1 \to t}'} \equiv 1  
    \}}
\newcommand{\Ckprime}[1]{\{
    \rho: \ts(\rho(i_k \to a_k), \rho)= \ts_k;\ts(\rho(i_k \to a_k)\circ \sigma_{k+1 \to j}, \rho(i_k \to a_k)\circ \sigma_{k+1 \to j}')= \ts_j, \forall #1+1 \leq j \leq t ;
    \rho(i_k) = \bar{a_k} ,
    F|_{\rho\circ(i_k \to a_k) \sigma_{k+1 \to t}} \equiv 0,
    F|_{\rho\circ(i_k \to a_k) \sigma_{k+1 \to t}'} \equiv 1 
    \}}
    
\newcommand{\calD}{\mathcal{D}}

\newcommand{\Dk}[1]{\{
    \rho: \ts(\rho(i_j \to a_j)\circ \sigma_{k+1 \to j}, \rho\circ \sigma_{k+1 \to j}')= \ts_j, \forall #1 \leq j \leq t ;
    \rho(i_k) = a_k ,
    F|_{\rho\circ \sigma_{k+1 \to t}} \equiv 0,
    F|_{\rho(i_j \to \bar{a_j})\circ \sigma_{k+1 \to t}'} \equiv 1  
    \}}
\newcommand{\Dkprime}[1]{\{
    \rho: \ts(\rho\circ \sigma_{k+1 \to j}, \rho(i_j \to \bar{a_j})\circ \sigma_{k+1 \to j}')= \ts_j, \forall #1+1 \leq j \leq t ;
    \rho(i_k) = a_k ,
    F|_{\rho\circ \sigma_{k+1 \to t}} \equiv 0,
    F|_{\rho(i_j \to \bar{a_j})\circ \sigma_{k+1 \to t}'} \equiv 1  
    \}}


where $\calA := \A$.

Define the map 
    $$\phi_a(\rho) = \rho \circ (\nvf(F,\rho) \to a)$$

    \begin{align*}
        \sum_{\rho \in \calA} Pr[\rho]
        &= \sum_{\rho \in \calA_t} \prob{}{\rho} \cdot \prod_{k=1}^{t-1} \frac{\sum_{\rho \in \calA_k}\prob{}{\rho}}{\sum_{\rho \in \calA_{k+1}}\prob{}{\rho}}\\
        &= \cdot \prod_{k=1}^{t-1} \frac{\sum_{\rho \in \calA_k}\prob{}{\rho}}{\sum_{\rho \in \calA_{k+1}}\prob{}{\rho}}\\
    \end{align*}

where 
$\calA_t = \left\{                        
            \rho: \nvf(F,\rho \circ \sigma_{t-1} = x_{i_t}), F|_{\rho \circ \sigma_t \equiv 0 , F|_{\rho \circ \sigma_t'} \equiv 1}    
            \right\}$   \\

$\calA_k = \left\{                        
            \rho: \nvf(F,\rho \circ \sigma_{j-1} = x_{i_j})\ \forall k \leq j \leq t , F|_{\rho \circ \sigma_t \equiv 0 , F|_{\rho \circ \sigma_t'} \equiv 1}    
            \right\}$\todo{stars}

  \begin{align*}
        \sum_{\rho \in \calA_t} \prob{}{\rho} 
        &= \sum_{\rho \in \calA_t} \sqrt{\prob{}{\rho \circ \sigma_t}}\sqrt{\prob{}{\rho \circ \sigma_t'}}  \\
        &\leq \sqrt{\sum_{\rho \in \calA_t}\prob{}{\rho \circ \sigma_t}}\sqrt{\sum_{\rho \in \calA_t}\prob{}{\rho \circ \sigma_t'}}\\
        &= \sqrt{\sum_{\rho \in \calB_t}\prob{}{\rho }}\sqrt{\sum_{\rho \in \calB_t'}\prob{}{\rho}}
  \end{align*}


Where $\calB_t:= \{
    \rho: \ts^{(t)}(\rho, \rho(i_t \to \bar{a_t}))= \ts_t,\ \rho(i_t) \to a_t,\ 
    F|_{\rho} \equiv 0,
    F|_{\rho (i_t \to \bar{a_t})} \equiv 1  
     \}$ and \\
     $\calB_t':= \{
    \rho: \ts^{(t)}(\rho(i_t \to a_t), \rho)= \ts_t,\ 
    \rho(i_t) = \bar{a_t},\     
    F|_{\rho(i_t \to a_t)} \equiv 0,
    F|_{\rho } \equiv 1  
     \}$.
Note that $\calB_t$ depends only on $x_{i_t}$.

    \begin{align*}
        \frac{\sum_{\rho \in \calA_k}\prob{}{\rho}}{\sum_{\rho \in \calA_{k+1}}\prob{}{\rho}} 
        &=\frac{\sum_{\rho \in \calA_k} \sqrt{\prob{}{\rho\circ \sigma_k}} \sqrt{\prob{}{\rho\circ \sigma_k'}}}{\sum_{\rho \in \calA_{k+1}}\prob{}{\rho\circ \sigma_k}}\\
        &= \sum_{\rho \in \calA_k} \sqrt{\frac{\prob{}{\rho\circ \sigma_k}}{\sum_{\rho \in \calA_{k+1}}\prob{}{\rho\circ \sigma_k}}} \sqrt{\frac{\prob{}{\rho\circ \sigma_k'}}{\sum_{\rho \in \calA_{k+1}}\prob{}{\rho\circ \sigma_k}}} \\
        &\leq \sqrt{\frac{\sum_{\rho \in \calA_k}\prob{}{\rho\circ \sigma_k}}  {\sum_{\rho \in \calA_{k+1}}\prob{}{\rho\circ \sigma_k}}} \sqrt{\frac{\sum_{\rho \in \calA_k}\prob{}{\rho\circ \sigma_k'}}{\sum_{\rho \in \calA_{k+1}}\prob{}{\rho\circ \sigma_k}}}\\        
    \end{align*}

%$\sum_{\rho \in \calA_{k+1}}\prob{}{\rho\circ \sigma_k} = \sum_{\rho \in \calB_k} \prob{}{\rho}$.

%
    % \ts^{(j)}(\rho \circ \sigma_{k+1 \to j}, \rho \circ \sigma_{k+1 \to j}')= \ts_j, \rho(i_j) = * \forall k+1 \leq j \leq t  ; 
    % F|_{\rho\circ \sigma_t} \equiv 0,
    % F|_{\rho \circ \sigma_t'} \equiv 1  
       
       
    % $\calC_{k}:= \left\{
    % \rho: \ts(\rho(i_k \to a_k),\rho(i_k \to a_k))= \ts_k;\ts(\rho(i_k \to a_k) \circ \sigma_{k+1 \to j}, \rho(i_k \to a_k) \circ \sigma_{k+1 \to j}')= \ts_j, \rho(i_j) = * \forall k+1 \leq j \leq t  ;
    % \rho(i_j) = a_j \forall 1 \leq j \leq k-1; \rho(i_k) = \bar{a_k};
    % F|_{\rho\circ \sigma_t} \equiv 0,
    % F|_{\rho \circ \sigma_t'} \equiv 1  
    %  \right\}$
    
    % $\calC_{k}':= \left\{
    % \rho: \ts(\rho(i_k \to a_k) \circ \sigma_{k+1 \to j}, \rho(i_k \to a_k) \circ \sigma_{k+1 \to j}')= \ts_j, \rho(i_j) = * \forall k+1 \leq j \leq t  ;
    % \rho(i_j) = a_j \forall 1 \leq j \leq k-1;\\
    % \rho(i_k) = \bar{a_k};
    % F|_{\rho\circ \sigma_t} \equiv 0,
    % F|_{\rho \circ \sigma_t'} \equiv 1  
    %  \right\}$
     
    \begin{align*}
        \frac{\sum_{\rho \in \calA_k}\prob{}{\rho\circ \sigma_k}}{\sum_{\rho \in \calA_{k+1}}\prob{}{\rho\circ \sigma_k}}
        &= \frac{\sum_{\rho \in \calB_k} \prob{}{\rho \circ (i_k \to a_k)}}{\sum_{\rho \in \calB_k} \prob{}{\rho \circ \nvf(F,\rho) \to a_k}}\\
        &= \frac{\sum_{\substack{\rho:\rho \in \calB_k\\ \ts(\rho, \rho(i_k \to \bar{a_k}))= \ts_k }} \prob{}{\rho}}{\sum_{\rho \in \calB_k} \prob{}{\rho}}
    \end{align*}
    Similarly,
    \begin{align*}
        \frac{\sum_{\rho \in \calA_k}\prob{}{\rho\circ \sigma_k'}}{\sum_{\rho \in \calA_{k+1}}\prob{}{\rho\circ \sigma_k'}}
        &= \frac{\sum_{\rho \in \calC_k} \prob{}{\rho \circ {i_k \to \bar{a_k}}}}{\sum_{\rho \in \calC_k'} \prob{}{\rho}}\\
        &= \frac{\sum_{\substack{\rho:\rho \in \calC_k'\\ \ts(\rho(i_k \to a_k), \rho)= \ts_k }} \prob{}{\rho}}{\sum_{\rho \in \calC_k'} \prob{}{\rho}}
    \end{align*}
        
        % %\newcommand{\nvfconds}{
    \{\rho| \nvf(F,\rho \circ \sigma_{k-1}) = x_{i_k} \forall 1 \leq k \leq t, F|_{\rho\circ \sigma_t} \equiv 0, F|_{\rho \circ \sigma_t'} \equiv 1  \}
}

\newcommand{\nvfcondsk}[1]{
\{\rho| \nvf(F,\rho \circ \sigma_{j-1}) = x_{i_j} \forall 1 \leq j \leq #1\}
}



\begin{lemma} 
    $$\prob{\rho}{\cdt(F,\rho) \mathrm{\ has\ depth\ } t} \leq p^t \sqrt{L(F)} {L(F)}^{t-1}$$ 
\end{lemma}


\begin{claim}
    $\nvf(F,\rho) = x_i \implies \prot(\rho\circ (x_i \to a), \rho\circ(x \to \bar{a})) = x_i$ and $\rho(i) = *$
\end{claim}


\begin{proof}
If $\cdt(F,\rho)$ has depth $t$, then there exist $a \in \cube{t}$ and variables $x_{i_1}, \ldots , x_{i_t}$ such that for all $k \in [t]$, $\nvf(F, \rho \circ \sigma_{k-1} )= x_{i_k}$,  
    where, $\sigma_k' = (x_{i_1}\to a_1,\ldots, x_{i_k} \to \bar{a_k})$ $\forall k \in [t]$ and $\sigma_0 = \sigma'_0 = ()$.

    More over, there exists a protocol $\prot$ such that if $\nvf(F, \rho\circ \sigma_{k-1}) = x_{i_k}$,  
    $\prot(\rho\circ \sigma_k, \sigma_k') = x_{i_k}$. Let $\calL_k$ be collection of transcripts which output $x_{i_k}$.   
    
    \begin{align*}
        \prob{\rho}{\cdt(F,\rho) \mathrm{\ has\ depth\ } t}   
        &\leq \sum_{x_{i_1}, \dots , x_{i_t}}  
        \prob{\rho}{\event{t}}    \\
        % &\leq \sum_{x_{i_1}, \dots , x_{i_t}}  
        % \prod_{k=1}^t \prob{\rho}{\event{k} \middle|
        % \event{k-1}}\\
        % &\leq \sum_{x_{i_1}, \dots , x_{i_t}} \sum_{\rho: (i),(ii) \forall 1\leq j \leq t,(iii)} \prob{}{\rho}
    \end{align*}
    

\end{proof}


Let $\calA = \nvfconds$

% Let $\calA_k:= \left\{
%     \rho: \ts(\rho\circ \sigma_j, \rho\circ \sigma_j')= \ts_j, 
%     \rho(i_j) = * \forall 1 \leq j \leq k ;
%     F|_{\rho\circ \sigma_k} \not\equiv \const
%     \right\}$
     
    \begin{align*}
        \sum_{x_{i_1}, \ldots ,x_{i_t}}\sum_{\rho \in \calA} Pr[\rho]
        &= \sum_{x_{i_1}, \ldots ,x_{i_t}} \prod_{k=1}^{t} \frac{\sum_{\rho \in \calA_{k}}\prob{}{\rho}}{\sum_{\rho \in \calA_{k-1}}\prob{}{\rho}}
    \end{align*}
where $\calA_k = \nvfcondsk{k}$ for each $k \in [t]$ and $\calA_0 = \{\rho| \rho: [n] \to \bra{0,1,*} \}$.
Note that $\calA_{k+1} = \{ \rho| \rho \in \calA_k , \nvf(F,\rho \circ \sigma_k) = x_{k+1} \}$. 
    \begin{align*}
        \sum_{x_{i_1}, \ldots ,x_{i_t}} \prod_{k=1}^{t} \frac{\sum_{\rho \in \calA_{k}}\prob{}{\rho}}{\sum_{\rho \in \calA_{k-1}}\prob{}{\rho}}
        &= \sum_{x_{i_1}} S_1(x_1) \sum_{x_{i_2}} S_2(x_{i_1},x_{i_2}) \ldots \sum_{x_{i_t}}S_t(x_{i_1},\ldots , x_{i_t})
    \end{align*}
    where $S_k(x_{i_1}, \ldots , x_{i_k}) = \frac{\sum_{\rho \in \calA_{k}}\prob{}{\rho}}{\sum_{\rho \in \calA_{k-1}}\prob{}{\rho}} $.
    
    Consider the expression $\sum_{x_{i_k}} S_k(x_{i_1}, \ldots , x_{i_k})$. Let $\sigma_k = \sigmaj{k}$ and $\sigma_k' = \sigmaprimej{k}$.


    \begin{align*}
        \sum_{x_{i_k}} \frac{\sum_{\rho \in \calA_{k}}\prob{}{\rho}}{\sum_{\rho \in \calA_{k-1}}\prob{}{\rho}}
        &=\cbra{\frac{2p}{1-p}}\sum_{x_{i_k}} \sum_{\rho \in \calA_k} \sqrt{\frac{\prob{}{\rho\circ \sigma_k}}{\sum_{\rho \in \calA_{k-1}}\prob{}{\rho\circ \sigma_{k-1}}}} \sqrt{\frac{\prob{}{\rho\circ \sigma_k'}}{\sum_{\rho \in \calA_{k-1}}\prob{}{\rho\circ \sigma_{k-1}}}}\\         
        &= \cbra{\frac{2p}{1-p}}\sum_{x_{i_k}} \sum_{ \substack{\rho \in \calB_{k-1} \\ \nvf(F,\rho) = x_{i_k}} }  \sqrt{\frac{\prob{}{\rho\circ (x_{i_k} \to a_k)}}{\sum_{\rho \in \calB_{k-1}}\prob{}{\rho}}} \sqrt{\frac{\prob{}{\rho\circ (x_{i_k} \to \bar{a_k})}}{\sum_{\rho \in \calB_{k-1}}\prob{}{\rho}}}\\
        &= \cbra{\frac{2p}{1-p}}\sum_{x_{i_k}} \sum_{ \substack{\rho \in \calB_{k-1} \\ \nvf(F,\rho) = x_{i_k}} }  \sqrt{\frac{\prob{}{\rho\circ (x_{i_k} \to a_k)}}{\sum_{\rho \in \calB_{k-1}}\prob{}{\rho\circ (x_{i_k} \to *)}+\prob{}{\rho\circ (x_{i_k} \to a_k)}+\prob{}{\rho\circ (x_{i_k} \to \bar{a_k})}}} \sqrt{\frac{\prob{}{\rho\circ (x_{i_k} \to \bar{a_k})}}{\sum_{\rho \in \calB_{k-1}}\prob{}{\rho}}}\\
    \end{align*}
    where $\calB_k = \{ \rho \circ \sigma_k | \rho \in \calA_k\}$.\\

\newcommand{\calX}{\mathcal{X}} 

    Consider the communication game $\game_{F,\calB_{k-1},a_k}$.
    The set of inputs for Alice $\calX_k = \{\rho \circ \}$
    \begin{align*}
        &\leq \sqrt{\frac{\prob{}{\rho\circ \sigma_k}}{\sum_{\rho \in \calA_{k-1}}\prob{}{\rho\circ \sigma_{k-1}}}} \sqrt{\frac{\prob{}{\rho\circ \sigma_k'}}{\sum_{\rho \in \calA_{k-1}}\prob{}{\rho\circ \sigma_{k-1}}}}
    \end{align*}

    
%\sum_{\substack{\rho: \rho \in \calB_{k-1}\\ \ts\cbra{\rho \circ (x_{i_k} \to a_k), \rho \circ(x_{i_k} \to \bar{a_k})} = \ts_k}}

    
    
    %Fix variables $x_{i_1}, \ldots , x_{i_t}$ and $a \in \bra{0,1}^t$.

%Define $\calB_{k} = \{\rho \circ (x \to a_k): \rho \in \calB_{k-1} \ \text{and}\ \nvf(F,\rho) = x\}$ and $\calB_0 = \calA_0$.  

%   \begin{align*}
%         \sum_{\rho \in \calA_1} \prob{}{\rho} 
%         &= \frac{2p}{1-p} \sum_{\rho \in \calA_1} \sqrt{\prob{}{\rho \circ \sigma_1}}\sqrt{\prob{}{\rho \circ \sigma_1'}}  \\
%         &\leq \sqrt{\sum_{\rho \in \calA_1}\prob{}{\rho \circ \sigma_1}}\sqrt{\sum_{\rho \in \calA_1}\prob{}{\rho \circ \sigma_1'}}        
%   \end{align*}


%$\calB_{k}:= \{
    % \rho: \ts(\rho, \rho(i_k \to \bar{a_k}))= \ts_k;
    % \ts(\rho(i_{j+1}= \ldots = i_k =*) , \rho(i_j \to \bar{a_j})(i_{j+1}= \ldots = i_k=*) )= \ts_j;
    % \rho(i_j) = a_j \forall 1 \leq j \leq k;
    % F|_\rho \not\equiv \const
    %  \}$
     
    % $\calB_{k}':= \{
    % \rho: \ts(\rho, \rho(i_k \to \bar{a_k}))= \ts_k;
    % \ts(\rho(i_{j+1}= \ldots = i_k =*) , \rho(i_j \to \bar{a_j})(i_{j+1}= \ldots = i_k=*) )= \ts_j;
    % \rho(i_j) = a_j \forall 1 \leq j \leq k;
    % F|_\rho \not\equiv \const
    %  \}$ 

    
    % $\calC_{k}:= \left\{
    % \rho: \ts(\rho(i_k \to a_k),\rho)= \ts_k;\ts(\rho(i_k \to a_k) \circ \sigma_{k+1 \to j}, \rho(i_k \to a_k) \circ \sigma_{k+1 \to j}')= \ts_j, \rho(i_j) = * \forall k+1 \leq j \leq t  ;
    % \rho(i_j) = a_j \forall 1 \leq j \leq k-1; \rho(i_k) = \bar{a_k};
    % F|_{\rho\circ \sigma_t} \equiv 0,
    % F|_{\rho \circ \sigma_t'} \equiv 1  
    %  \right\}$
    
    % $\calC_{k}':= \left\{
    % \rho: \ts(\rho(i_k \to a_k) \circ \sigma_{k+1 \to j}, \rho(i_k \to a_k) \circ \sigma_{k+1 \to j}')= \ts_j, \rho(i_j) = * \forall k+1 \leq j \leq t  ;
    % \rho(i_j) = a_j \forall 1 \leq j \leq k-1;\\
    % \rho(i_k) = \bar{a_k};
    % F|_{\rho\circ \sigma_t} \equiv 0,
    % F|_{\rho \circ \sigma_t'} \equiv 1  
    %  \right\}$
    
         
    % \begin{align*}
    %     \frac{\sum_{\rho \in \calA_k}\prob{}{\rho\circ \sigma_k}}{\sum_{\rho \in \calA_{k+1}}\prob{}{\rho\circ \sigma_k}}
    %     &= \frac{\sum_{\rho \in \calB_k} \prob{}{\rho}}{\sum_{\rho \in \calB_k'} \prob{}{\rho}}\\
    %     &= \frac{\sum_{\substack{\rho:\rho \in \calB_k'\\ \ts(\rho, \rho(i_k \to \bar{a_k}))= \ts_k }} \prob{}{\rho}}{\sum_{\rho \in \calB_k'} \prob{}{\rho}}
    % \end{align*}
    % Similarly,
    % \begin{align*}
    %     \frac{\sum_{\rho \in \calA_k}\prob{}{\rho\circ \sigma_k'}}{\sum_{\rho \in \calA_{k+1}}\prob{}{\rho\circ \sigma_k'}}
    %     &= \frac{\sum_{\rho \in \calC_k} \prob{}{\rho}}{\sum_{\rho \in \calC_k'} \prob{}{\rho}}\\
    %     &= \frac{\sum_{\substack{\rho:\rho \in \calC_k'\\ \ts(\rho(i_k \to a_k), \rho)= \ts_k }} \prob{}{\rho}}{\sum_{\rho \in \calC_k'} \prob{}{\rho}}
    % \end{align*}
    
    % Let,\\
    % $cond_j(\rho): \condj{j}$ and $\starsj{j}$\\
    % $cond_\top(\rho) : F|_{\rho \circ \sigma_t} \equiv 1$\\
    % $cond_\bot(\rho) : F|_{\rho \circ \sigma_t'} \equiv 0$.    
    
    
    % Let $\mathcal{L}_j = set\ of\ leaves\ of\ CC_F^{(j)}$\\
    % $\mathcal{L}_j(i_j)= subset\ of \mathcal{L}_j \ labelled\ by\ x_{i_j}$.
    
    % $S_1:= \sum_{\substack{\rho:cond_1(\rho) \\ F|_{\rho_1} \neq         const}}\sqrt{Pr[\rho_1]\cdot Pr[\rho_1']}$
    
    % $S_j:= \frac{\sum_{\substack{\rho:cond_j(\rho):1 \leq j \leq         i\\F|_{\rho_j} \neq const}}\sqrt{Pr[\rho_i]\cdot             Pr[\rho_i']}}{\sum_{cond_j(\rho):1\leq j \leq i-            1}Pr[\rho_{i-1}]} \ \forall 2 \leq j \leq t$
    
    % Observe that $S_i$ doesn't depend on $x_{i_j},\ell_j$ for $j >i$.
    
    % $\sum_{x_{i_1},\dots ,x_{i_t}}\sum_{\substack{\ell_j \in \mathcal{L}_1(x_{i_1})}\\ \vdots \\ \ell_t \in \mathcal{L}_t(x_{i_t}} \prod_{i=1}^t S(i) = \sum_{\substack{x_{i_1}\\ \ell_1 \in \mathcal{L}_1(x_{i_1})}}S(1) \dots \sum_{\substack{x_{i_t}\\ \ell_t \in \mathcal{L}_t(x_{i_t})}}S_t$        
            
        
           

% \newcommand{\defrk}{
%     \{
%         \rho \circ \sigma_{k-1}: \prot^{(j)}(\rho \circ \sigma_j,\rho \circ \sigma_j') = \ts_j \forall 1 \leq j \leq k-1           
%     \}
% }

%     Let $X^{(t)} = \{ \phi_{a_t}(\rho): F|\phi_{a_t}(\rho) \equiv 0, F|\phi_{\bar{a_t}}(\rho) \equiv 1\}$
%     $Y^{(t)} = \{ \phi_{\bar{a_t}}(\rho): F|\phi_{a_t}(\rho) \equiv 0, F|\phi_{\bar{a_t}}(\rho) \equiv 1 \}$.
    
%     $X^{(k)} = \{ 
%         \phi_{a_k}(\rho): 
%             \prot^{(k)}(\phi_{a_k}(\rho),\phi_{\bar{a_k}}(\rho)) = \ts_k,
%             \prot^{(j)}(\phi_{a_k}(\rho)\circ \sigma_{k+1 \to j}, \phi_{a_j}(\rho) \circ \sigma_{k+1 \to j}') = \ts_j,
%             F|_{\phi_{a_k}(\rho)\circ\sigma_{k+1 \to t}} \equiv 0,
%             F|_{\phi_{a_k}(\rho)\circ\sigma_{k+1 \to t}'} \equiv 1
%     \}
%             $
            
    
%     Let $\calR_k = \defrk$ and $\calR_1 = \{\rho: \nvf(F,\rho) \neq \bot \}$. 
%     Consider the games $\calG_k$ for each $k \in [t]$ where $\calG_k = \game_{F,R_k,a_k}$.    
% \tulasi{is $R_k$ this or from below? }
%     By Lemma \ref{lem: protocol}, there exists a protocol $\prot^{(k)}$ for the game $\calG_k$ for each $k \in [t]$ whose number of transcripts is $|\prot^{(k)}|$.      
%     Note that $\calG_k$, $\prot^{(k)}$ depend only the variables $x_{i_1}, \ldots, x_{i_{k-1}}$.     